\documentclass[../veristore_security_model.tex]{subfiles}

\begin{document}
\section{Entry Points}

% COMMENT: List all interfaces where external entities interact with veri-store. Each entry point
% represents a potential attack surface. Group related operations into logical categories rather
% than listing every API method. For each entry point, note its trust level and whether it sends
% or receives data (or both). Examples: client API for store/retrieve, inter-server verification
% protocol, storage subsystem file operations.

\subsection{Client API}

The client API is the main way that users interact with veri-store. It groups three high-level operations:
\begin{itemize}
    \item \textbf{Store operation}: The client sends a data block to veri-store. The client encodes the block into fragments, generates the fpcc, and then issues HTTP requests to store each fragment on its assigned server.
    \item \textbf{Retrieve operation}: The client requests a stored block by its identifier. It asks multiple servers for fragments, verifies each fragment against the fpcc, and reconstructs the original data once enough valid fragments are available.
    \item \textbf{Verify operation}: The client (or an operator) can re-check that stored fragments are still consistent with the fpcc, without needing to reconstruct the full block. This detects tampering after the initial store.
\end{itemize}

\subsection{Inter-Server Communication}

In the basic project deployment, all ``servers'' run on the same host, but the design assumes they could be separate machines on a network. Inter-server communication is limited and well-defined:
\begin{itemize}
    \item \textbf{Fingerprint exchange}: In extended deployments, servers may exchange fingerprint values or fpcc digests to support cross-checking and audit workflows. This path is treated as semi-trusted and is protected by the same fingerprinting assumptions as the client API.
    \item \textbf{Fragment retrieval}: A server may request fragments from peer servers to assist with recovery or migration. Any fragment obtained this way is still subject to the same fpcc-based verification before being used.
\end{itemize}

\subsection{Storage Subsystem}

The storage subsystem covers all local file operations used to persist fragments and metadata on each server:
\begin{itemize}
    \item \textbf{Fragment write}: The server writes encoded fragments and their associated metadata (including fpcc digests and verification status) to disk. Writes use a write-then-rename pattern to reduce the chance of partial files.
    \item \textbf{Fragment read}: The server reads fragments and metadata back from disk to answer client requests or to perform re-verification. If a file is missing or malformed, the server treats the fragment as unavailable.
    \item \textbf{Fingerprint storage}: The server stores enough information (such as fpcc digests or JSON-serialized fpcc) to allow future integrity checks without needing the client to resend the fpcc.
\end{itemize}

\newpage
\end{document}